\section{Insertion into a regime classifier}\label{sec:insertion}

A measured reliability is of use only where it can enter a decision, and we test
the two points at which it can enter a regime classifier built on a rolling
two-state Gaussian hidden Markov model, in the specification of \citet{blake2025}. Neither
insertion works, and the two failures share a common cause that has nothing to do
with the size of $\lambda$.

\subsection{Shrinkage of the observable}

The natural insertion is to shrink the observable toward its unconditional mean
with weight $\lambda$, replacing $\log RV_M$ with
$(1-\lambda)\mu + \lambda \log RV_M$ before classification. The specification
standardises the observable within each rolling window, and for any window $W$
the shrunk series has mean $(1-\lambda)\mu + \lambda\,\overline{x}_W$ and
standard deviation $\lambda\, s_W$, so the relation
\begin{equation}\label{eq:affine}
\frac{z - \overline{z}_W}{s_W(z)} \;=\; \frac{x - \overline{x}_W}{s_W(x)}
\end{equation}
holds identically for every $\lambda > 0$, leaving the standardised input
unchanged and the classification unable to move. Across all eight cells the two
arms differ in $0$ states, and the reduction in misclassification is $0.00$ % numbers.csv: k11_red_vs_A1[ES/GLOBEX/1day]
percentage points in every cell, in sample and out. % numbers.csv: k11_red_vs_A1[NQ/RTH/30min]

Equation~\eqref{eq:affine} is not specific to reliability shrinkage but holds for
any affine transformation of the observable, so the entire class of linear proxy
corrections applied upstream of a standardising classifier is annihilated, not
attenuated. The recoverable band is empty, not small, and the conclusion is
available analytically, before any classification has been run.

\subsection{Noise in the observation equation}

The remaining route is to hand the model the noise instead of the correction,
decomposing the emission variance in each state into a state variance plus a
known observation-noise variance $v = (1-\lambda)\operatorname{Var}(\log RV_M)$
and holding $v$ fixed while the rest is estimated. On synthetic data with known
state variance and known added noise the estimator behaves as intended, in that
the free-variance model recovers the total emission variance and the fixed-noise
model recovers the state variance by subtracting $v$, with the state means of the
two agreeing to ten decimal places.

That correctness is what defeats the exercise. With
$\sigma_k^2 = \max(\text{total}_k - v, 0)$ the emission variance is
$\max(\text{total}_k, v)$, so the fixed-noise model is the free-variance model
equipped with a variance floor at $v$, identical to it wherever both states clear
that floor. The floor is far from idle, since at the headline scaling it binds in
$46{,}120$ of $72{,}619$ windows, or $63.5$ percent, and moves $18{,}311$ % numbers.csv: k12_windows_binding_s100_ext, k12_windows_total_s100_ext, k12_binding_share_s100_ext
in-sample states. Even so, the maximum reduction in misclassification is $0.00$ % numbers.csv: k12_states_differ_insample_s100_ext
percentage points in every cell at every scaling, and where the floor binds % numbers.csv: k12_max_reduction_pp
hardest the classification degrades.

The reason is that a variance floor cannot sharpen a discrimination. Where both
states carry the same known uncertainty, knowing how uncertain an observation is
tells a two-state model nothing, because the common term enters both emission
densities and cancels in the posterior ratio. Only the floor survives that
cancellation, and it widens both densities and pulls marginal observations toward
the nearer state mean.

\subsection{Scope of the negative results}

Both results are specific to a discrete-state model with state-independent noise,
and a model whose latent variable is continuous does not share the cancellation.
In a
stochastic volatility state-space model with a log-variance signal equation and a
known observation-error variance, the noise enters a filtering recursion and not
a two-way likelihood ratio. We have not tested such a model and make no claim
about its behaviour.
